% !TeX root = ../../Book3.tex
\section{Phragmén–Lindelöf 原理}
\label{b3:sec:1904}

最大值原理在无界扇形中不能直接应用。Phragmén--Lindelöf 方法先乘一个衰减指数，在大圆弧压住增长，再把有限区域的最大值原理送到无穷远。

\subsection{扇形域}
\label{b3:subsec:190401}

扇形 \(S=\{r e^{\mathrm{i}\theta}:r>0,\ |\theta|<\beta/2\}\)，\(0<\beta<2\pi\)，具有可选定的幂分支 \(z^p\)。若 \(p\beta/2<\pi/2\)，则 \(\operatorname{Re}z^p\ge c|z|^p\)，其中 \(c=\cos(p\beta/2)>0\)。这给整个闭扇形同方向的衰减，不能只检查中心射线。

\subsection{增长条件}
\label{b3:subsec:190402}

\term[phragmen-lindelof]{Phragmén--Lindelöf 原理}[Phragmén--Lindelöf principle]的一个常用版本如下。
\begin{theorem}[Phragmén--Lindelöf：扇形]\label{b3:thm:19-pl}
设 \(f\) 在上述扇形全纯并连续到有限边界与顶点。若两边界射线上 \(|f|\le1\)，且 \(|f(z)|\le C e^{A|z|^q}\)，其中 \(q<\pi/\beta\)，则 \(|f|\le1\) 于整个扇形。
\end{theorem}
\begin{proof}
选 \(q<p<\pi/\beta\)，令 \(f_\varepsilon=f e^{-\varepsilon z^p}\)。两射线上模不超过一；大圆弧上模不超过 \(C e^{Ar^q-\varepsilon c r^p}\)，充分大时亦不超过一。顶点值由连续性不超过一。对截断扇形用最大模原理得 \(|f_\varepsilon|\le1\)，再令大半径趋无穷，最后 \(\varepsilon\downarrow0\)。
\end{proof}
临界指数必须另设条件：\(e^{z^{\pi/\beta}}\) 在两边界模一，在中心射线却发散，给出严格不等式的必要性。

\subsection{整函数应用}
\label{b3:subsec:190403}

整函数若阶 \(\rho<1\)，并在实轴有界，则为常数。选择 \(\rho<q<1\)，阶的定义给 \(Ce^{Ar^q}\) 增长；分别对上、下半平面（角宽 \(\pi\)）应用扇形原理，得到全平面有界，再用 Liouville。

阶一的 \(e^{\mathrm{i}z}\) 在实轴模一而非常值，显示阈值不能只把严格不等号改成不大于。若函数仅在一个半平面全纯，结论只保证内部有界，不会自动成为常数。

\subsection{唯一性定理}
\label{b3:subsec:190404}

同一方法给唯一性。若 \(f,g\) 在扇形满足相同的次临界增长，且在两边界相等，则差 \(h\) 的边界零。对每个 \(\varepsilon>0\)，缩放 \(h/\varepsilon\) 满足边界模不超过一，扇形原理给 \(|h|\le\varepsilon\)，故 \(f=g\)。

这里边界相等不能直接当作恒等定理的内部聚点条件。本证明通过增长界得到全域比较；若边界为可作 Schwarz 反射的直线，还能用反射后的恒等定理得到更强的局部唯一性。临界例子中的零边值实部 \(\operatorname{Re}z^{\pi/\beta}\) 也说明无穷远会引入额外自由度。


\ProseParagraphBreak

扇形的旋转不改变临界阈值。若中轴辐角为 \(\theta_0\)，应在衰减因子中使用 \((e^{-\mathrm{i}\theta_0}z)^p\)，使实部在整个扇形有正下界；直接保留 \(z^p\) 可能在旋转后变成增长因子。幂的分支固定在该扇形中，连续延到两边界射线；顶点只有连续值，不必声称非整数幂在顶点全纯。截断扇形应用最大模原理时，顶点属于边界，内域不包含那个可能的分支点。

\ProseParagraphBreak

若 \(f\) 在顶点仅有界而未声明连续，证明可以先删去小圆。小圆上的模上界可能大于一，不能直接丢弃；若顶点为可去奇点，补上后再使用原定理即可。若允许顶点有极点，则边界射线在接近顶点仍受模一控制已经强烈限制它；但对于只控制远处边界的变式，小圆会留下额外常数。定理的有限边界与顶点假设用来控制截断域的全部边界，不能只关注大圆弧的增长。

\ProseParagraphBreak

一个多射线应用可以把阶阈值提高。若整函数阶 \(\rho\) 小于 \(\pi/\beta\)，并在若干从原点发出的射线上有共同模界 \(M\)，且相邻射线最大角宽为 \(\beta\)，则每个夹角扇形都可选增长指数 \(q\) 位于 \(\rho\) 与 \(\pi/\beta\) 之间。分别使用扇形定理，得到全平面 \(|f|\le M\)，故常数。例如四条坐标半轴把平面分成角宽 \(\pi/2\) 的扇形，于是阶小于二且在实轴、虚轴都有界 的整函数为常数。临界阶二的 \(e^{z^2}\) 在虚轴有界，在实轴却无界；\(e^{\mathrm{i}z^2}\) 在两轴均模一而非常值，给出正确的临界反例。

\ProseParagraphBreak

这里有限多个扇形的常数可以统一，因为只取有限次最大值；若射线无限多而各自界不同，不应未经估计就宣称有一个全平面共同界。把每条射线上的界与一个统一界区别开，也是使用 Liouville 结论前必须核查的量词。

\ProseParagraphBreak

边界零的唯一性应用需要同时假设两函数的增长。差函数的界由两者相加获得，常数可能加倍但指数仍是同一个次临界 \(q\)。然后对任意 \(\delta>0\) 使用 \(h/\delta\)，增长常数会随 \(\delta\) 改变；定理允许每次使用不同有限常数，因此最终 \(|h|\le\delta\) 合法。不能要求所有 \(\delta\) 共用一个增长常数，那会把证明中的无害缩放误认为失败。

\ProseParagraphBreak

无界区域比较还有一个常用的条带版本，它让边界模界不同的情形呈现出插值形式。
\begin{theorem}[有界函数的三线不等式]\label{b3:thm:19-three-lines}
设 \(f\) 在 \(0<\operatorname{Re}z<1\) 全纯，在闭条带连续有界，并满足 \(|f(\mathrm{i}t)|\le M_0\)、\(|f(1+\mathrm{i}t)|\le M_1\)，其中 \(M_0,M_1>0\)。则
\[
 |f(x+\mathrm{i}y)|\le M_0^{1-x}M_1^x\qquad(0\le x\le1).
\]
\end{theorem}
\begin{proof}
令 \(g(z)=f(z)\exp(-(1-z)\log M_0-z\log M_1)\)。实对数无分支问题，\(g\) 在闭条带有界，两竖边模不超过一。对 \(\varepsilon>0\)，令 \(g_\varepsilon(z)=g(z)e^{\varepsilon z^2-\varepsilon}\)。在两竖边有 \(\operatorname{Re}(\varepsilon z^2-\varepsilon)=\varepsilon(x^2-y^2-1)\le0\)。在横边 \(y=\pm T\)，其模不超过 \(\|g\|_\infty e^{-\varepsilon T^2}\)，充分大 \(T\) 时也不超过一。矩形最大模原理给 \(|g_\varepsilon|\le1\)，令 \(T\to\infty\)、\(\varepsilon\downarrow0\) 得 \(|g|\le1\)，还原即得结论。
\end{proof}
有界性在这条证明中控制横边，指数修正把横边送到无穷远后消去。若希望允许更快的竖向增长，必须选择新的修正并重新估计，不能直接引用当前版本。三线不等式的权重由观察点在两边界之间的相对位置决定，正如三圆定理的权重由对数半径决定；两个公式背后都是求出合适的调和比较函数。
\ProseParagraphBreak
如果某一边界模界为零，可先把该 \(M_j\) 替换为任意小正数，应用已经证明的正界版本再取极限；于是整个条带函数为零。这一处理不需要对零取对数，也不会留下未定义的指数因子。一般的边界参数极限应在不等式已经建立之后执行，而不是在构造辅助全纯函数时直接代入奇异参数。
